\section{Round-by-round and knowledge soundness}\label{sec:rbr}

\subsection{Hypothesis and standing assumptions}\label{sec:hyp}

Thresholds and degree bounds are those of Annex~B (\S\ref{sec:annexb}): $\dgr^{(i)}=2^{\kappa_i}-1+\ell_i$ and $\agr^{(i,j)}=\agr_i-(\ell_i-j)$ for one free integer $\agr_i\le\blen_i$ per level. The point error of level $i$ is $\epp_i:=(\dgr^{(i)}+2^{\ell_i}-1)/\q$, as in \eqref{eq:eps}.

\begin{hypothesis}[$\Lb_i(L_i)$]\label{hyp:LB}
For every $0\le m\le\ell_i$ and every interleaved word $\orc:\cube{m}\times\dom_i\to\E$: $|\lst(\orc;\agr^{(i,0)},\dgr^{(i)})|\le L_i$.
\end{hypothesis}
\status{A hypothesis, not a claim of this paper. It is \stpr{} in the two cases named below (the Johnson bound and the counting bound) and open otherwise; \S\ref{sec:conj} states what is \stcj.}

Lists at a larger agreement are subsets, so the same $L_i$ bounds every list the proof uses. In fact the proof uses the hypothesis only on a diagonal: for a word with $r$ lane variables, $0\le r\le\ell_i$, at the agreement $\agr^{(i,\ell_i-r)}$ (the committed word at $\agr^{(i,0)}$, a word folded $k$ times at $\agr^{(i,k)}$, a single lane at $\agr_i$). This weaker form is the exact hypothesis of the machine-checked composition (\S\ref{sec:formal}); it differs from the stated one by at most $\ell_i$ columns and changes no number in this paper. We state the stronger form because it is the one the reductions of \S\ref{sec:interleave} deliver. The hypothesis is a statement about the code alone. It is a theorem in two cases: with the Johnson count $L_i=\jl(\agr^{(i,0)},\dgr^{(i)})$ when $(\agr^{(i,0)})^2>\blen_i\dgr^{(i)}$ (Theorem~\ref{thm:johnson}), and with the trivial \emph{counting bound} $L_i=\binom{\blen_i}{\dgr^{(i)}+1}/\binom{\agr^{(i,0)}}{\dgr^{(i)}+1}$ for $\agr^{(i,0)}\ge\dgr^{(i)}+1$ (each tuple is determined by any $\dgr^{(i)}+1$ of its agreement columns), which is astronomically large except at the last, smallest levels. Section~\ref{sec:lists} reduces it to single-lane words over $\K$ and says what is known beyond these two cases.

\paragraph{Standing assumptions.} All are satisfied by every parameter set leanVM derives ($\ell_0=6$, $\ell_i=4$ after, $\kappa_{\nlev-1}\in\{2,\dots,5\}$, $\blen_i\le2^{28}$).
\begin{enumerate}[itemsep=1pt,topsep=2pt,label=\textnormal{(S\arabic*)}]
\item $1\le\ell_i\le2^{\kappa_i}$, so that $\Wn{\kappa_i}H^{(i)}$ with $\deg H^{(i)}<\ell_i$ is the part of index $2^{\kappa_i}$ and above of a polynomial of degree at most $\dgr^{(i)}$.
\item $\dgr^{(i)}+2\le\agr^{(i,0)}$ and $\agr_i\le\blen_i$.
\item $\binom{\blen_i}{2}\le2\q$.
\item $\Kz$ is an integer with $\Kz>M:=2\q$, and $0<\dsc<1$. These two are constants of the analysis; they do not appear in the protocol.
\end{enumerate}
We write $\thr:=\Kz(\Kz-1)(\Kz-2\q)$. The machine-checked composition shows which of these assumptions the soundness statement really uses: $\ell_i\ge1$, a nonempty domain, (S3), $M\ge2\q$, $\dsc>0$, $\thr>0$, thresholds that grow by at least one per round, and $\Lb$. It uses neither (S2), nor $\dsc<1$, nor $\agr_i\le\blen_i$. (S2) is needed only for the uniqueness of the near tuple in the descent of \S\ref{sec:ks}, that is, to read the extraction predicate as the output of a decoder, and (S1) enters only through the two properties of the high part used in the transition M3.

\subsection{The theorem}\label{sec:thm}

Call the input claims (the weighted claims $(W_k,c_k)_k$ and the commit-phase point claims $(z_m,y_m)_m$ of level $0$) \emph{false} if no $U\in\lst(\orc^{(0)};\agr^{(0,0)},\dgr^{(0)})$ has $\ev^{(0)}_{z_m}(U)=y_m$ for all $m$ and $\inner{W_k}{g_U}=c_k$ for all $k$.

\begin{theorem}[Round-by-round soundness of the DEEP-fold protocol]\label{thm:rbr-deep}
Assume \textnormal{(S1)} to \textnormal{(S4)} and $\Lb_i(L_i)$ for every level. Then Protocol~\ref{prot:D} is round-by-round sound with the following errors.
\begin{center}\upshape
\renewcommand{\arraystretch}{1.9}
\begin{tabular}{@{}lp{0.3\textwidth}p{0.56\textwidth}@{}}
\toprule
 & verifier message & round-by-round error\\
\midrule
M0 & first points of level $i$ & $\displaystyle\binom{L_i}{2}\epp_i^{\,\bar s^{(i)}_1}+\frac{L_i^3\,\q^6}{\thr}\sum_{j=1}^{\ell_i}\dsc^{-(j-1)}\,\epp_i^{\,\bar s^{(i)}_j}$\\
M1 & batching $\lambda_0$ & $(J_0-1)/\q$; with the sixteen ring-switching scalars in the same message, $(J_0+1)/\q$\\
M2 & fold challenge of round $j$ (with batch $j+1$) & $(\Kz-1)/\q^2+\dsc$ for $j<\ell_i$; $(\Kz-1)/\q^2$ for $j=\ell_i$\\
M3 & queries of level $i<\nlev-1$, with $\lambda_{i+1}$ & $\bigl((\agr_i-1)/\blen_i\bigr)^{\nq_i}+\bigl(\bar s^{(i)}_{\ell_i}+\nq_i\bigr)/\q$\\
M4 & queries of the last level & $\bigl((\agr_{\nlev-1}-1)/\blen_{\nlev-1}\bigr)^{\nq_{\nlev-1}}$\\
M5 & each tail round & $2/\q$\\
\bottomrule
\end{tabular}
\end{center}
For one batch of $s_i$ points the first row is $\bigl[\binom{L_i}{2}+\frac{L_i^3\q^6}{\thr}\sum_{j<\ell_i}\dsc^{-j}\bigr]\epp_i^{\,s_i}$.
\end{theorem}
\status{\stpr{} (this section), and \stmc{} in the following form. Each transition is a standalone counting statement with exactly this error: M0 is \lean{dens_not_good_le} (one batch) and \lean{dens_not_goodS_le} (a schedule) in \leanfile{Protocol/Level.lean} and \leanfile{Protocol/Ladder.lean}; M1 \lean{batch_transition_density}; M2 \lean{fold_transition_density}, together with \lean{levelPot_step} or \lean{ladderPot_step} of the same two \texttt{Protocol} files for the term $\dsc$; M3 \lean{query_transition_density}; M4 \lean{final_transition_density}; M5 \lean{tail_round_transition_density}; the five transition statements are in \leanfile{Rounds/}. For one batch of points per level the composition is machine-checked for an abstract model of the protocol: \lean{protocolD_rbr} in \leanfile{Soundness/ProtocolD.lean} states that the model, on any ladder of levels with any tail, has a round-by-round state function (\lean{StateFunction}, properties (i) to (iv), in \leanfile{Soundness/RBR.lean}) with exactly the errors of rows M0 to M5, and \lean{protocolD_sound}, \lean{protocolD_sound_plain} in \leanfile{Soundness/Sound.lean} that every prover strategy wins with probability at most their sum. In that state function property (i) holds at every start, with no condition on the claims, and the falsity of the claims is part of the winning condition, as explained in \S\ref{sec:state}; the potential is the one-batch potential of the Lean development (status of Lemma~\ref{lem:T}). In the model a committed word is a function sent by the prover, the map $U\mapsto\mpoly_{g_U}$ is any linear map, the high part is any polynomial with the two properties used in M3, the tail is a sumcheck on two abstract tables, row M1 is $(J_0-1)/\q$, and the errors are those of the interactive protocol before any grind. That the model has prover strategies is machine-checked on every ladder and with every tail, under two conditions on the abstract maps without which the prover can have no message (\lean{strategy_nonempty_ladder} in \leanfile{Soundness/Sound.lean}; \S\ref{sec:composed}). No Lean theorem covers: the composition for a schedule with later batches (only its point message), the identities of Annex~B that make the real encoder an instance of the model, completeness, the Fiat-Shamir transformation and grinding, Merkle commitments, or the outer protocol (\S\ref{sec:formal}).}

Neither the Johnson hypothesis nor the correlated agreement error $\epsaff$ of Theorem~\ref{thm:annexb} appears, and no list size multiplies a term of order $1/\q$: the list size enters only the first row, where it is multiplied by $\epp_i^{s}$ and can be paid for with points. With $\Kz=2^9\q$ and $\dsc=2^{-119}$ the fold message has error below $2^{-118}$ and $2^{-125}$ after its $7$-bit grind; the query message is as in Annex~B; the first row is brought below $2^{-110}$ (so $2^{-125}$ after the $15$-bit grind) by the choice of the schedule (\S\ref{sec:points}); the first opening message has error at most $513/\q$ and needs a $7$-bit grind where Annex~B needs $13$.

\begin{corollary}[At the Johnson threshold the protocol is unconditionally sound]\label{cor:johnson}
If $(\agr^{(i,0)})^2>\blen_i\dgr^{(i)}$ for every level, Theorem~\ref{thm:rbr-deep} holds with $L_i=\jl(\agr^{(i,0)},\dgr^{(i)})$ and no hypothesis.
\end{corollary}
\status{\stpr{} (Theorem~\ref{thm:rbr-deep} and Theorem~\ref{thm:johnson}). The Lean development does not connect the Johnson bound to the composed theorem; its instances without a list hypothesis are in the unique-decoding range $\agr^{(i,0)}>(\blen_i+\dgr^{(i)})/2$, in \leanfile{Soundness/Hypothesis.lean}: \lean{protocolD_sound_unique}, and two concrete instances of it in which every hypothesis is discharged: \lean{protocolD_sound_concrete_unique_le}, one level with no tail over a field of $2^{128}$ elements on which every prover strategy wins with probability at most $2^{-109}$, and \lean{protocolD_sound_concrete_ladder}, two levels with a tail of two rounds over the same field. Neither is a deployed parameter set (\S\ref{sec:composed}): they show that the hypotheses of the composed theorem can be met.}

This is the precise sense in which the DEEP fold changes what has to be proved. The corollary changes the proof and not the cost: at these thresholds the query counts are those of Annex~B up to one query per level, and the points cost about what that saves (\S\ref{sec:sizes}). Below the Johnson radius the theorem is conditional on $\Lb$, and \emph{no instance of $\Lb$ below the Johnson radius with a useful $L_i$ is known} (\S\ref{sec:lists}), apart from the counting bound at the smallest levels.

\subsection{The state function}\label{sec:state}

Let $\orc^{(i,j)}$ be the level-$i$ word folded by its first $j$ challenges, so $\orc^{(i,0)}=\orc^{(i)}$ and $\orc^{(i,\ell_i)}=w^{(i)}$, and let $\lst^{(i,j)}:=\lst(\orc^{(i,j)};\agr^{(i,j)},\dgr^{(i)})$, of size at most $L_i$. Let $\pot^{(i)}_k$ be the potential \eqref{eq:pot} of level $i$, with the exceptional sets taken at $\agr=\agr^{(i,j)}$, $\dgr=\dgr^{(i)}$, $L=L_i$, the theorem's $\Kz$ and $M=2\q$, which is $\max(2\q,\binom{\blen_i}{2})$ by (S3). The first points $\mathbf z^{(i)}$ of level $i$ are \emph{good} for $\orc^{(i)}$ if they are separating for $\orc^{(i)}$ at $(\agr^{(i,0)},\dgr^{(i)})$ (Lemma~\ref{lem:P}) and $\pot^{(i)}_0<1$.

The invariant is: some check of $\bV$ has already failed, or the condition attached to the last verifier message holds.
\begin{itemize}[itemsep=1pt,topsep=2pt]
\item After the first points of level $0$, and through the outer protocol up to the opening: $\mathbf z^{(0)}$ is good for $\orc^{(0)}$. At the start of the opening, in addition: the input claims are false.
\item After $\lambda_i$ ($j=0$; for $i\ge1$ this is the query message of level $i-1$) and after the fold message of round $j\ge1$:
\begin{itemize}[itemsep=0pt,topsep=1pt]
\item[$(\mathrm I_{i,j})$] no $U\in\lst^{(i,j)}$ satisfies the running claim $(W^{(i,j)},\rc{i,j})$ and all carried point claims $(z_m,v_m)$ of the batches $1,\dots,\max(j,1)$ at round $j$;
\item[$(\mathrm G_{i,j})$] $\pot^{(i)}_j<1$.
\end{itemize}
\item After the first points of level $i+1$: $(\mathrm I_{i,\ell_i})$ holds and $\mathbf z^{(i+1)}$ is good for $\orc^{(i+1)}$.
\item After the last query message: the low part $f^{(\nlev)}$ of the polynomial sent does not satisfy the running claim.
\item After each tail round: the running sumcheck claim is false.
\end{itemize}

\paragraph{Properties (i), (ii), (iv).} (i) When the commit-phase answers and the claims are given and false, the start is the definition of false input claims together with goodness, whose failure is charged to M0. (ii) Every condition is a function of the transcript up to the last verifier message only. Goodness and $\pot^{(i)}_j$ depend on the commitment, the points and the challenges. $(\mathrm I_{i,j})$ depends in addition on $\rc{i,j}$ and the carried values, which are functions of prover messages sent before the challenge of round $j$. The conditions after a point message do not involve the answers. So no prover message breaks the invariant. (iv) At the end either a check failed or the last sumcheck claim is false, and $\bV$, who holds $f^{(\nlev)}$, evaluates it and rejects.

\paragraph{Commitment and opening as one protocol.} The commit-phase answers and the input claims are sent by the prover \emph{after} the level-$0$ points. So when commitment and opening are read as one protocol, ``the input claims are false'' cannot be part of a state that survives prover messages: the prover leaves it by naming true claims. The statement that composes is the relativised one of Theorem~\ref{thm:K} below with the canonical predicate $\Bad(U):\Longleftrightarrow U$ violates a commit-phase point claim or an input claim. With it the state before the points is empty (the standing assumptions and $\Lb$), the hypothesis of the transition M1 is a tautology, and property (iv) reads: at a doomed end position, if $\bV$ has accepted then the final polynomial descends to a member of $\lst^{(0,0)}$ that matches all commit-phase answers and satisfies all input claims. Theorem~\ref{thm:rbr-deep} is the corollary: if no such member exists, $\bV$ rejects. This is the form of the machine-checked composition, and it came out of the formalization. One consequence is worth stating: the errors of M1 and of all later messages are paid on every run, whether the claims are true or false, which is what a statement about adaptively chosen claims needs.

\subsection{The transitions}\label{sec:transitions}

\paragraph{M0, first points of level $i$.} $\orc^{(i)}$ is committed before its points. By Lemma~\ref{lem:P} with the $\bar s^{(i)}_1$ points of batch $1$, they are separating except with probability $\binom{L_i}{2}\epp_i^{\bar s_1}$. A fold of a word of the shape of round $j$ has $\ell_i-j$ lane variables, and its list at agreement $\agr^{(i,j)}\ge\agr^{(i,0)}$ is contained in its list at $\agr^{(i,0)}$, which has at most $L_i$ members by $\Lb_i$; so Lemma~\ref{lem:T}(a) applies with $L=L_i$ and $\pot^{(i)}_0\ge1$ with probability at most $(L_i^3\q^6/\thr)\sum_j\dsc^{-(j-1)}\epp_i^{\bar s_j}$. For $i\ge1$ the condition $(\mathrm I_{i-1,\ell_{i-1}})$ does not involve the new points and stays.

\paragraph{M1, batching $\lambda_0$.} The points are separating, so at most one $U^\ast\in\lst^{(0,0)}$ satisfies all point claims $(z_m,y_m)$. If there is none, $(\mathrm I_{0,0})$ holds for every $\lambda_0$. Otherwise the input claims are false, so $U^\ast$ violates some weighted claim, and $\sum_k\lambda_0^{k-1}(\inner{W_k}{U^\ast}-c_k)$ is a nonzero polynomial of degree at most $J_0-1$ in $\lambda_0$: it vanishes with probability at most $(J_0-1)/\q$. Every other list member violates a point claim, which is not batched, so \emph{there is no union over the list}. $(\mathrm G_{0,0})$ is the goodness of $\mathbf z^{(0)}$, which $\lambda_0$ does not touch.

\paragraph{M2, fold message of round $j$.} Before it, $(\mathrm I_{i,j-1})$ and $(\mathrm G_{i,j-1})$ hold, and the prover has sent the answers to batch $j$ (if any) and its $4+\bar s_j$ scalars. By Lemma~\ref{lem:T}(b) the $\bar s_j$ points held are not in $\Exc_j(\orc^{(i,j-1)})$. Apply Lemma~\ref{lem:F}(2) with $\orc=\orc^{(i,j-1)}$, $\agr=\agr^{(i,j)}$, $(W,c^\circ)=(W^{(i,j-1)},\rc{i,j-1})$, the carried values $v_m$ of all points held, and $B$ true. Its hypothesis is implied by $(\mathrm I_{i,j-1})$: $\agr-1=\agr^{(i,j-1)}$, and the answers to batch $j$ add point claims, which can only shrink the set of list members satisfying all of them. Its conclusion is about $P\in\lst(\orc_c;\agr^{(i,j)},\dgr^{(i)})=\lst^{(i,j)}$ with (p), which says that $P$ satisfies the updated carried claims \eqref{eq:carry}, and (w), which says that $P$ satisfies the updated running claim \eqref{eq:claim-update}. So $(\mathrm I_{i,j})$ fails for fewer than $\Kz$ of the $\q^2$ challenges, whatever the fresh points of the same message are. By Lemma~\ref{lem:T}(b), $(\mathrm G_{i,j})$ fails with probability less than $\dsc$, and never for $j=\ell_i$, where $\pot_{\ell_i}=0$.

\paragraph{M3, query message of level $i<\nlev-1$.} Before it the prover has committed $\orc^{(i+1)}$, sent $\phigh^{(i)}$, received the points $\mathbf z^{(i+1)}$ and answered them; $(\mathrm I_{i,\ell_i})$ holds and $\mathbf z^{(i+1)}$ is good. By separation at most one $U^\ast\in\lst^{(i+1,0)}$ satisfies all new point claims, and it is fixed before the queries. If there is none, $(\mathrm I_{i+1,0})$ holds. Otherwise let $g^\ast:=g_{U^\ast}$ and $P^\ast:=\mpoly_{g^\ast}+\Wn{\kappa_i}H^{(i)}$, a polynomial of degree at most $\dgr^{(i)}$ by (S1), with low part $g^\ast$, fixed before the queries.
\begin{itemize}[itemsep=1pt,topsep=2pt]
\item \emph{Case 1: $P^\ast$ agrees with $w^{(i)}=\orc^{(i,\ell_i)}$ on at least $\agr_i=\agr^{(i,\ell_i)}$ columns.} Then $P^\ast\in\lst^{(i,\ell_i)}$, and by $(\mathrm I_{i,\ell_i})$ either $P^\ast$ violates the running claim, and then $g^\ast$ violates the residual claim (a weighted claim sees only the low part); or $P^\ast(\zeta_m)\ne v_m$ for some carried claim (at round $\ell_i$ a point claim is a plain evaluation), and then $\inner{W_{\zeta_m}}{g^\ast}=\mpoly_{g^\ast}(\zeta_m)=P^\ast(\zeta_m)+\Wn{\kappa_i}(\zeta_m)H^{(i)}(\zeta_m)\ne c'_{i,m}$: $g^\ast$ violates a converted point claim \eqref{eq:convert}.
\item \emph{Case 2: otherwise.} All $\nq_i$ independent uniform queries land on agreement columns with probability at most $((\agr_i-1)/\blen_i)^{\nq_i}$. If some query has $P^\ast(x)\ne w^{(i)}(x)$, adding $\Wn{\kappa_i}(x)H^{(i)}(x)$ to both sides gives $\inner{W_x}{g^\ast}=\mpoly_{g^\ast}(x)\ne c_{i,k}$: $g^\ast$ violates a query claim.
\end{itemize}
So except with that probability $U^\ast$ violates one of the $J_{i+1}$ claims of the level boundary, which are fixed by the queries; $\lambda_{i+1}$ is uniform and independent of them, the claims carry distinct powers $0,\dots,\bar s^{(i)}_{\ell_i}+\nq_i$, and the batch stays violated except with probability $(\bar s^{(i)}_{\ell_i}+\nq_i)/\q$. Every other member of $\lst^{(i+1,0)}$ violates a new point claim. This is $(\mathrm I_{i+1,0})$, and $(\mathrm G_{i+1,0})$ is the goodness of $\mathbf z^{(i+1)}$.

\paragraph{M4, last query message.} The prover has sent $P:=P^{(\nlev-1,\ell)}$, of degree at most $\dgr^{(\nlev-1)}$ by its coefficient count, with low part $f^{(\nlev)}$. If $P$ agrees with $w^{(\nlev-1)}$ on at least $\agr_{\nlev-1}$ columns, then $P\in\lst^{(\nlev-1,\ell)}$ and by $(\mathrm I_{\nlev-1,\ell})$ either $P(\zeta_m)\ne v_m$ for some $m$, and the point check of the last level has failed, or $P$, hence $f^{(\nlev)}$, violates the running claim. Otherwise some query has $P(x)\ne w^{(\nlev-1)}(x)$ except with probability $((\agr_{\nlev-1}-1)/\blen_{\nlev-1})^{\nq_{\nlev-1}}$, and the query check fails.

\paragraph{M5, tail round.} A degree-two sumcheck round from a false claim leaves a false claim except with probability $2/\q$ \cite{LFKN92}. This completes the proof of Theorem~\ref{thm:rbr-deep}. \qed

\paragraph{Boundary cases.} The first round of a level applies Lemma~\ref{lem:F} to the committed word itself, whose list $\lst(\orc;\agr^{(i,1)}-1,\dgr)=\lst^{(i,0)}$ is the list Lemma~\ref{lem:P} pins. At the last round the fold is one lane, $\ev^{(\ell)}$ is evaluation at $\zeta$, and there is no $\dsc$. All lists of a level are at the one degree $\dgr^{(i)}$: Lemma~\ref{lem:L2} returns halves of degrees $(\dgr-1,\dgr)$ from folds of degree $\dgr$, so nothing grows going back, and the point claims evaluate whole polynomials, so they pin the slack coefficients too. The high-part correction is used twice in M3, and $\phigh^{(i)}$ is sent before the points and the queries. The point check of the last level is necessary: without it a polynomial of the last list that satisfies the running claim and violates a point claim would be accepted. Points that fall in $\dom$, equal points, or $\theta=0$ need no special treatment: all are inside the Schwartz-Zippel counts.

\paragraph{Level $0$.} The lanes of $\orc^{(0)}$ are $\K$-valued and committed before any outer claim exists. A member of $\lst^{(0,0)}$ agrees with $\K$-valued lanes on at least $\agr^{(0,0)}\ge\dgr+1$ columns of $\dom_0\subseteq\K$, so its coefficients are in $\K$ (Lagrange interpolation on $\dgr+1$ of those columns has nodes and values in $\K$). Goodness of $\mathbf z^{(0)}$ is a property of $(\orc^{(0)},\mathbf z^{(0)})$ and survives every message of an outer protocol. The claims arrive after the points and may depend on them and on the answers, which is harmless because Lemma~\ref{lem:F}(2) holds for every weight, claim and carried values at fixed non-exceptional points. The folds $\orc^{(0,j)}$, $j\ge1$, are $\E$-valued, which is why $\Lb_0$ is stated for $\E$-valued lanes.

\subsection{Knowledge soundness}\label{sec:ks}

As in the existing analysis of Annex~B, soundness and extraction are one statement relative to a predicate: ``every \emph{bad} candidate violates the claims'' is preserved by every verifier message, and an accepting run that preserved it yields, by a descent, a candidate that is not bad.

\begin{definition}[Near tuples, pull-back]\label{def:pull}
A \emph{near tuple} of a word $\orc$ on a set $S\subseteq\dom$ is a tuple of degree at most $\dgr$ agreeing with $\orc$ on all but at most one column of $S$. For a round $j\ge1$ of level $i$, the word $\orc:=\orc^{(i,j-1)}$, the challenge $c$ of round $j$ and a predicate $B$ on tuples of round $j-1$, the \emph{pull-back} of $B$ is the predicate on tuples $P$ of round $j$
\[
(\pull B)(P)\;:\Longleftrightarrow\;\Bigl[\,|\agrset(P,\orc_c)|\ge\agr^{(i,j)}\ \Longrightarrow\ \text{every near tuple of }\orc\text{ on }\agrset(P,\orc_c)\text{ satisfies }B\,\Bigr],
\]
where $\agrset$ is the set of agreement columns. Given a predicate $\Bad$ on level-$0$ tuples, let $B_{0,0}:=\Bad$, $B_{i,j}:=\pull B_{i,j-1}$ along the challenge of round $j$, and across levels $B_{i+1,0}(U):\Longleftrightarrow B_{i,\ell_i}(P^\ast_U)$ with $P^\ast_U:=\mpoly_{g_U}+\Wn{\kappa_i}H^{(i)}$. The relativised condition $(\mathrm I^B_{i,j})$ is: no $U\in\lst^{(i,j)}$ \emph{with $B_{i,j}(U)$} satisfies the running claim and all carried point claims.
\end{definition}

\begin{lemma}[Uniqueness of near tuples]\label{lem:near}
If $|S|\ge\dgr+3$, a word has at most one near tuple on $S$.
\end{lemma}
\status{\stmc: \lean{nearT_unique} in \leanfile{Rounds/Extract.lean}.}
\begin{proof}
Two near tuples agree with the word, hence with each other, on all but at most two columns of $S$: at least $\dgr+1$ columns, more than their degree.
\end{proof}

\begin{theorem}[Round-by-round knowledge soundness]\label{thm:K}
Assume the hypotheses of Theorem~\ref{thm:rbr-deep}. Let $\Bad$ be a predicate on level-$0$ tuples, fixed before the batching challenge $\lambda_0$ (it may depend on the commit-phase answers and on the claims), such that every $U\in\lst^{(0,0)}$ with $\Bad(U)$ that satisfies the commit-phase point claims violates one of the input claims. Then every verifier message preserves the relativised invariant except with the error of Theorem~\ref{thm:rbr-deep} for that message, and on every accepting run that preserved it the final polynomial $P$ lies in $\lst^{(\nlev-1,\ell)}$ and does not satisfy $B_{\nlev-1,\ell}$. From such a $P$ the descent below computes a tuple $U_0\in\lst(\orc^{(0)};\agr^{(0,0)},\dgr^{(0)})$ with $\neg\Bad(U_0)$.
\end{theorem}
\status{\stpr{} (below). \stmc{} in the following form: the relativised transitions for an arbitrary predicate, \lean{batch_transition}, \lean{fold_transition}, \lean{query_transition}, \lean{final_transition} in \leanfile{Rounds/}; one step of the descent with uniqueness, \lean{descent_step} in \leanfile{Rounds/Extract.lean}; and, for the canonical predicate of \S\ref{sec:state} and one batch of points per level, the composition over a run in the abstract model: \lean{protocolD_sound_extract} in \leanfile{Soundness/Sound.lean} bounds by the sum of the errors the probability that $\bV$ accepts and the final polynomial does not descend to a member of $\lst^{(0,0)}$ matching the answers and satisfying the claims, and \lean{descend_iff_unique}, \lean{descend_step} in \leanfile{Soundness/Extract.lean} show, under (S2) and for one fold round, that the extraction predicate (which asks for \emph{some} near tuple) is the statement about \emph{the} unique near tuple; that over a whole run it is the chain of unique near tuples follows by unfolding its definition and is not a Lean theorem. Not machine-checked: the composition for a general predicate or a schedule with later batches, the extractor as an algorithm (that the near tuple is found by one-error decoding, and in polynomial time), and the extraction of the committed words from the adversary's hash queries.}
\begin{proof}
The transitions of \S\ref{sec:transitions} with these changes. M0 is unchanged: goodness does not involve $B$. M1: at most one $U^\ast$ satisfies the point claims; if there is none or $\neg\Bad(U^\ast)$, $(\mathrm I^B_{0,0})$ holds, and otherwise $U^\ast$ violates an input claim by hypothesis, with the same count. M2: Lemma~\ref{lem:F}(2) with $B:=B_{i,j-1}$; its hypothesis is $(\mathrm I^B_{i,j-1})$, and its clause (b) for a $P\in\lst^{(i,j)}$ is exactly $B_{i,j}(P)$, because the premise of the pull-back holds for a list member; so fewer than $\Kz$ challenges break $(\mathrm I^B_{i,j})$. M3: let $U^\ast$ be the member of $\lst^{(i+1,0)}$ satisfying the new point claims; if $\neg B_{i+1,0}(U^\ast)$, $(\mathrm I^B_{i+1,0})$ holds, and otherwise $B_{i,\ell_i}(P^\ast)$ and the two cases apply with $(\mathrm I^B_{i,\ell_i})$ in Case~1, with the same counts. M4: if $P\notin\lst^{(\nlev-1,\ell)}$ a query fails except with the stated probability; if $P$ is in the list and $B_{\nlev-1,\ell}(P)$, then by $(\mathrm I^B_{\nlev-1,\ell})$ the point check fails or $f^{(\nlev)}$ violates the running claim, and the tail rejects except with M5's errors. So an accepting run that escaped no message has $P\in\lst^{(\nlev-1,\ell)}$ and $\neg B_{\nlev-1,\ell}(P)$.

\emph{Descent.} $\neg(\pull B)(P)$ says: $P$ agrees with the folded word on at least $\agr^{(i,j)}$ columns, and some near tuple $T$ of the unfolded word on that agreement set has $\neg B(T)$. Such a $T$ has degree at most $\dgr$ and agrees with the unfolded word on at least $\agr^{(i,j)}-1=\agr^{(i,j-1)}$ columns, so $T\in\lst^{(i,j-1)}$. By (S2), $\agr^{(i,j)}\ge\dgr^{(i)}+3$ for $j\ge1$, so by Lemma~\ref{lem:near} it is \emph{the} near tuple, and it fails $B$. Going back through the rounds of the last level gives $T\in\lst^{(\nlev-1,0)}$ with $\neg B_{\nlev-1,0}(T)$, that is $\neg B_{\nlev-2,\ell}(P^\ast_T)$, which again contains ``$P^\ast_T$ is in the list of the previous level'', and so on down to level $0$.
\end{proof}

\paragraph{The extractor.} Input: the accepted transcript (challenges, high parts, final polynomial) and the committed words $\orc^{(i)}$, read off the adversary's hash queries as in the existing extractor of leanVM. Set $T\leftarrow$ the final polynomial. For $i=\nlev-1$ down to $0$ and $j=\ell_i$ down to $1$: compute $S:=\agrset(T,\orc^{(i,j)})$; fail if $|S|<\agr^{(i,j)}$; set $T\leftarrow$ the near tuple of $\orc^{(i,j-1)}$ on $S$, by decoding at most one error on the known set $S$; fail if none exists. After the rounds of level $i>0$, set $T\leftarrow\mpoly_{g_T}+\Wn{\kappa_{i-1}}H^{(i-1)}$. Output the low parts of the lanes of $T$.

The extractor \emph{never list-decodes}: every step is unique decoding of one tuple with at most one error on a set it knows. Below the Johnson radius no efficient list decoder for these codes is known and the lists may be enormous; the extractor needs neither, because it is led to its candidate by the final polynomial, which the accepting prover sent in the clear. It performs $\sum_i\ell_i$ decodings, is deterministic and straight-line (no rewinding), and its running time is polynomial in the size of the committed words, independent of every list size. The Lean development adds a fact the hand proof does not state: the near tuple folds back to the tuple it was decoded from (\lean{NearT.fold_eq}), so the chain of tuples is consistent with the challenges of the transcript.

\begin{corollary}[The extractor returns the pinned candidate]\label{cor:K2}
Suppose $\mathbf z^{(0)}$ is good, and let $U^\ast$ be the member of $\lst^{(0,0)}$ consistent with the commit-phase answers, if there is one. For every opening of every set of input claims, except with the errors of Theorem~\ref{thm:rbr-deep}: if the run accepts, then $U^\ast$ exists, $g_{U^\ast}$ satisfies every input claim, and the descent outputs $U^\ast$.
\end{corollary}
\status{\stpr. \stmc{} pieces: the predicate below satisfies the hypothesis of Theorem~\ref{thm:K} for any claims when the pinned member exists (\lean{k2_hfalse} in \leanfile{Rounds/Batch.lean}), and for good points at most one member of $\lst^{(0,0)}$ matches the answers and satisfies the claims (\lean{claimsHold_unique} in \leanfile{Soundness/Extract.lean}).}
\begin{proof}
Theorem~\ref{thm:K} with $\Bad(U):\Longleftrightarrow U\ne U^\ast$ or $g_U$ violates an input claim (everything is bad if there is no $U^\ast$). A bad list member satisfying the point claims is $U^\ast$ by separation, so it violates an input claim: the hypothesis holds. The descent's output is not bad.
\end{proof}

So the commitment is binding in the sense an outer protocol needs: the candidate is fixed once the prover has answered the commit-phase points, before any opening claim exists, and every later opening is about that candidate. One fact is inside the errors and worth recording: an accepted run need not have a pinned tuple. If a few members share their low part and the prover is off in one answer, the run can be accepted with probability about (number of such members)$/\q$ per fold attempt, below the fold error; it then opens true claims about the common low part, and the descent returns one of those tuples.

\subsection{The interface with leanVM's outer protocol}\label{sec:interface}

In leanVM the opening is the last phase of a larger protocol (bus, table sumchecks, ring switching), and today that protocol is analysed as follows: after the Merkle root, the commitment's OOD point leaves at most one \emph{candidate} stack answering the OOD claim, and every outer phase is a statement about that one $\K$-valued candidate, fixed before the first outer challenge. Protocol~\ref{prot:D} provides the same object with one substitution: ``the candidate answering the OOD claim'' becomes ``the low part of the list tuple answering the $\bar s^{(0)}_1$ commit-phase point claims'', which is unique by Lemma~\ref{lem:P}, is over $\K$ (\S\ref{sec:transitions}, level $0$), and is fixed before the first outer challenge. The OOD claim leaves the pool of weighted claims and the point claims are handed to the opening as its carried claims. The sixteen ring-switching scalars share a message with $\lambda_0$, and the ring switch is sound for one fixed $\K$-valued stack with error $2/\q$, which gives the $(J_0+1)/\q$ of row M1.

\status{This subsection is \stpr{} only at the level of the statements at the interface (which predicate singles the candidate out, and what the first opening message is charged). The outer phases themselves are inherited from leanVM's existing analysis and were not re-derived; only the ring switch was re-read against its statement. Adopting the protocol would be a change to both provers, the three verifiers and the formal development of leanVM.}
